% !TeX root = ../main.tex
% Base: borrador de JulianaEmanuel, commit 85e7c67. Se conservan sus políticas,
% la etiqueta original y la referencia clabAbadiPACELC2012; se amplía el análisis.
% Asistencia del borrador base: OpenAI Codex. Ampliación: OpenAI GPT-6.1 Sol.
\section{Análisis CAP y PACELC}
\label{sec:clab-cap}
\label{sec:clab-cap-pacelc}

CAP expresa que, bajo una partición, no se pueden garantizar simultáneamente
consistencia linealizable y respuesta satisfactoria para todas las solicitudes
dirigidas a nodos no fallidos\cite{clabGilbertLynchCAP}. Su disponibilidad no es
el porcentaje mensual de disponibilidad de un requisito no funcional. Tampoco
la consistencia de CAP equivale simplemente a las restricciones de integridad
de ACID. PACELC añade el compromiso, cuando no hay partición, entre latencia y
consistencia\cite{clabAbadiPACELC2012,clabAbadiPACELC}.

CLAB todavía no tiene replicación de negocio en el esqueleto de E1. Por eso no
se etiqueta PostgreSQL ni todo el sistema como ``CP'' o ``AP'' por su nombre.
Se especifican políticas objetivo por operación: \textbf{PC/EC} prefiere
consistencia bajo partición y también en operación normal; \textbf{PA/EL} permite
lecturas de una proyección local potencialmente atrasada y prioriza la latencia.
La segunda política solo mantiene lectura en los nodos que tengan una copia
accesible; no promete disponibilidad universal ni escritura multimaestro.

\subsection{Relación con los requisitos E1}

Se emplean los identificadores actuales de \texttt{main.tex} y
\texttt{Carlos.tex}, revisados en \texttt{b212d70}. RF-P02 exige como máximo
una confirmación ante dos solicitudes simultáneas para el mismo laboratorio
y horario. RNF-03 exige impedir cambios si no hay permiso verificable, si la
sesión está revocada o si Identidad no puede verificarla. Ambos orientan la
preferencia por consistencia en el camino de confirmación.

RNF-01 propone al menos el 95\,\% de consultas válidas de disponibilidad
correctas en menos de 2 segundos durante 10 minutos con 50 usuarios concurrentes.
RNF-02 propone al menos el 99\,\% de comprobaciones correctas dentro de 5 segundos,
con una sonda por minuto durante dos semanas de horario académico. Son metas
de evaluación, no valores medidos ni garantías deducidas de CAP. RNF-06 y
RNF-07 requieren confirmar y persistir las reservas de prueba con Avisos detenido
y mantener auditoría fuente de todos los cambios exitosos, respectivamente.

\subsection{Reservas y bloqueos de agenda: PC/EC}

\textbf{Dato e invariante.} El servicio Reservas es la única autoridad de
ocupación: dos reservas confirmadas no pueden solaparse para el mismo laboratorio;
un bloqueo confirmado impide nuevas reservas dentro de su intervalo. Un bloqueo
que afecta reservas existentes debe registrar las afectaciones en la misma
transacción. El catálogo de Laboratorios no es una segunda agenda.

\textbf{Con partición.} Si el solicitante no llega a la autoridad de agenda o
esta no puede confirmar sus dependencias de autorización y catálogo, no recibe
una confirmación nueva. No se aceptan reservas sin conexión para fusionarlas
después: el usuario debe reintentar y puede perder el horario. Un calendario
almacenado previamente, si se muestra, es informativo y no concede el intervalo.

\textbf{Sin partición.} La confirmación espera la transacción autoritativa.
Se propone serializar por laboratorio todas las reservas, bloqueos, cancelaciones
y recurrencias mediante un bloqueo de fila estable. Intervalos contiguos
se interpretan como $[inicio,fin)$: se permite que uno termine cuando comienza
el siguiente. Se sacrifica latencia frente a confirmar desde una copia atrasada.
La garantía debe probarse con concurrencia, no inferirse de usar
\texttt{@Transactional}.

\textbf{Recuperación.} Un timeout no prueba que la transacción haya fallado.
Se propone una clave de idempotencia por solicitud, registrada con la operación,
y una consulta del resultado antes de reintentar. El diseño de API aún debe
implementarse; el endpoint de comprobación de E1 no crea reservas.

\subsection{Usuarios activos, roles y permisos: PC/EC}

Identidad es la fuente de verdad de credenciales y autorizaciones. Los otros
servicios consultan su validación antes de aceptar peticiones protegidas.
Durante una partición con Identidad se rechaza temporalmente la operación;
no se presume que un token aún vigente mantenga permisos ya revocados.
Sin partición se paga una llamada adicional para verificar usuario activo y
permisos. Sigue pendiente acordar el efecto de una revocación que coincida con
una operación ya autorizada; no se promete anular retroactivamente una
transacción confirmada. Esta dependencia reduce disponibilidad y debe incluirse
en las pruebas y en la estimación del 99\,\% de disponibilidad.

\subsection{Catálogo, equipos e incidencias: autoridad PC/EC}

Laboratorios y Recursos decide actualizaciones de catálogo, estado operativo e
incidencias. Si no se llega a su almacén autoritativo no se confirman cambios
ni cierres de incidencias. Sin partición se confirma después de persistir el
cambio y su auditoría local. Esto impide que dos copias aisladas comuniquen
estados contradictorios como definitivos. La pérdida de este servicio no borra
reservas existentes, pero bloquea nuevas confirmaciones que necesiten validar
el catálogo. El cierre por mantenimiento requiere además un bloqueo de agenda
confirmado en Reservas; actualizar solamente un campo del catálogo no garantiza
el cierre del intervalo. No se promete atomicidad entre ambos servicios.

Para RF-P04 se propone que Laboratorios persista cierre pendiente y solicite
por I13 una barrera de admisión en Reservas. Solo publica fuera de servicio
después de confirmar durablemente esa barrera. Cada confirmación de reserva
la comprueba bajo el mismo bloqueo local, aunque una consulta previa haya leído
un catálogo activo. Una reserva confirmada antes de la barrera se conserva;
una confirmación posterior se rechaza. No se cancelan automáticamente reservas
existentes ni se autoriza con ello el uso físico de un laboratorio cerrado.

Durante una partición, el trámite no se anuncia completado sin confirmación.
Si la barrera se activó y se perdió la respuesta, el catálogo puede seguir
pendiente mientras Reservas rechaza nuevas confirmaciones. Se prefiere perder
disponibilidad antes que reabrir por una compensación insegura. La recuperación
consulta o reintenta la misma solicitud durable; identificador y revisión
evitan duplicados y órdenes anteriores aplicadas fuera de secuencia.

La reapertura comienza cambiando el catálogo a apertura pendiente, todavía
no reservable, y después retira la barrera operativa mediante I13. Solo tras
confirmar esa revisión publica activo. No retira otros bloqueos de agenda.
Una respuesta perdida conserva la apertura pendiente, no un retorno automático
a fuera de servicio con una barrera que pudiera estar retirada. Sin partición
se acepta la latencia de estas dos confirmaciones locales; no se promete una
transacción distribuida atómica. El protocolo completo queda en ADR-001 y aún
debe implementarse y validarse con el equipo.

\subsection{Asistencia y datos académicos: autoridad PC/EC}

Reservas registra asistencia únicamente para una reserva válida; Identidad y
Académico mantiene estudiantes, matrículas, periodos y asignaturas. Las escrituras
requieren acceso al servicio propietario y autorización vigente. Durante una
partición no se confirma una asistencia o matrícula solo desde el navegador.
Sin partición se prioriza la persistencia autoritativa sobre responder desde
una copia local. Una reserva puede conservar una instantánea de sus metadatos
académicos para explicar el registro histórico, sin convertirla en fuente de
verdad de matrículas o permisos.

\subsection{Auditoría: fuente PC/EC y vista central PA/EL}

La operación de negocio y su auditoría fuente se guardan en una transacción
local: si no puede registrarse esa auditoría obligatoria, no se confirma el
cambio. La copia central en Información y Avisos se alimenta posteriormente.
Durante una partición entre servicios, la copia central puede seguir mostrando
los eventos ya recibidos, con su marca de última actualización; no certifica
que incluya toda la actividad actual. Sin partición se prefiere leer esa
proyección sin consultar todos los orígenes en cada pantalla. La convergencia
requiere recuperación de registros durables y deduplicación, no entrega
``exactamente una vez'' supuesta.

\subsection{Reportes y notificaciones: proyecciones PA/EL}

Información y Avisos puede consultar su copia local aunque no lleguen cambios
de Reservas o Laboratorios. El usuario verá estadísticas atrasadas o recibirá
la notificación después; ninguna de esas vistas confirma disponibilidad de
horarios. En operación normal se conserva la lectura local para reducir latencia.
El atraso no tiene una cota demostrada: depende de la frecuencia de consulta,
el volumen pendiente y las interrupciones. Tras restablecer la conexión, se
reanuda desde cursores persistidos y se deduplican los cambios. Si el propio
servicio, su almacén o Identidad están inaccesibles, la consulta protegida no
queda disponible por llamarla PA/EL.

\subsection{Alcance de las garantías y comprobaciones pendientes}

Un único escritor por dominio simplifica la autoridad, pero no es una solución
de alta disponibilidad. Si en E3 se añaden réplicas para datos críticos, deberán
definirse replicación, confirmación durable y elección de primario con protección
contra dos escritores simultáneos; una promoción manual sin verificar ese
aislamiento no demuestra consistencia. Las políticas anteriores orientan ese
diseño y no acreditan un protocolo de consenso implementado. La replicación
debe preservar el orden de confirmación incluso al cambiar el nodo responsable.

La validación prevista debe comprobar: dos solicitudes concurrentes para el
mismo intervalo; reserva simultánea con bloqueo; pérdida de respuesta tras una
confirmación; desactivación concurrente de RF-P04; revocación de permisos;
corte y recuperación de los canales de
proyección; reinicio del consumidor conservando cursores; y fallo de la
persistencia de auditoría. Para los intervalos iguales, el criterio propuesto
es como máximo una confirmación y un resultado explícito para la otra solicitud.
Para la proyección, el criterio es convergencia sin pérdida ni duplicación
observable después de recuperar los cambios conservados en el origen.

\paragraph{Estado.}
Este análisis amplía el borrador incorporado por JulianaEmanuel en el commit
\texttt{85e7c67}, preparado con OpenAI Codex, conservando sus políticas de agenda,
autorización y proyecciones y su referencia de PACELC. La ampliación se preparó
con asistencia de IA (OpenAI, GPT-6.1 Sol) y se contrasta con ADR-001 y los
contenedores objetivo. No se ejecutaron
experimentos de partición, concurrencia ni replicación del sistema de negocio
durante su preparación. Sus efectos deben contrastarse con requisitos trazables
al levantamiento real y con la implementación posterior.
